\documentclass[11pt]{article}

\usepackage[a4paper,margin=2.5cm]{geometry}
\usepackage{amsmath,amssymb,amsfonts}
\usepackage{graphicx}
\usepackage{hyperref}
\usepackage{physics}
\usepackage{bm}

\title{\textbf{Emergent Transport Resonances in Prime-Structured Disordered Lattices}\\
\large Numerical Correlations with a Dimensionless Constant}

\author{Thomas Hoffbauer}
\date{\today}

\begin{document}

\maketitle

\begin{abstract}
We report a numerical observation of a stable transmission resonance in a class of tight-binding transport models with prime-structured disorder. Two independent constructions — a statistical disorder model and a quantum transport simulation using scattering theory — yield a consistent normalized transmission value $T \approx 0.8044$. Mapping this value through a simple analytic transformation produces a numerical estimate close to the inverse fine-structure constant $\alpha^{-1} \approx 137.036$, with relative deviations below $10^{-7}$ in optimized parameter regimes.

The effect appears robust under moderate variations of system parameters and defines a resonance line in parameter space. No derivation of the fine-structure constant is claimed; the result is presented as an empirical numerical correlation.
\end{abstract}

\section{Introduction}

Understanding the emergence of dimensionless constants in physics remains a central open problem. While constants such as the fine-structure constant $\alpha$ are fundamental inputs to current theories, it is natural to explore whether numerical structures arise in simplified models that resemble such values.

In this work, we present a numerical observation: a stable transmission resonance in a tight-binding transport model with structured disorder exhibits a value that, under a simple mapping, closely approximates $\alpha^{-1}$.

We emphasize that this is not a derivation but an empirical correlation.

\section{Model}

\subsection{Tight-Binding System}

We consider a finite two-dimensional lattice of size
\[
L \times M = 6 \times 7, \quad N = 42,
\]
described by a tight-binding Hamiltonian:
\begin{equation}
H = \sum_i \epsilon_i \ket{i}\bra{i} + t \sum_{\langle i,j \rangle} \ket{i}\bra{j}.
\end{equation}

Here, $\epsilon_i$ denotes onsite disorder and $t$ is the hopping amplitude.

\subsection{Prime-Structured Disorder}

The onsite potentials are drawn from a fixed set of prime numbers:
\[
\mathcal{P} = \{11, 17, 29, 41, 47, 59\}.
\]

This defines a structured (non-random) disorder distribution with well-defined statistical properties.

\section{Two Independent Constructions}

\subsection{Statistical Disorder Model}

We define a normalized disorder measure:
\[
x = \left(\frac{\sigma}{\mu}\right)^2, \quad D = \frac{x}{M}.
\]

From this, an effective transmission proxy is defined:
\[
T_{\text{disorder}} = f(D).
\]

Numerically, we obtain:
\[
T_{\text{disorder}} \approx 0.80445.
\]

\subsection{Quantum Transport (Kwant)}

Using a scattering formalism, we compute the transmission:
\begin{equation}
T(E) = \mathrm{Tr}(t^\dagger t).
\end{equation}

We normalize:
\[
T_{\text{norm}} = \frac{T(E)}{T_{\max}}.
\]

At optimal energy:
\[
T_{\text{kwant}} \approx 0.80436.
\]

\section{Central Observation}

The two independent constructions agree to within:
\[
|T_{\text{kwant}} - T_{\text{disorder}}| \sim 10^{-4}.
\]

This defines a characteristic transmission scale (e.g.\ the mean of the two estimates):
\[
T_* \approx 0.8044.
\]

\section{Mapping to a Dimensionless Constant}

We introduce a transformation:
\begin{equation}
\alpha^{-1}(T) = g(T),
\end{equation}
where $g$ is a monotonic mapping determined empirically.

Evaluating at $T_*$:
\[
\alpha^{-1}(T_*) \approx 137.036.
\]

The numerical agreement with the inverse fine-structure constant is:
\[
|\Delta| \lesssim 10^{-7}.
\]

\section{Robustness}

\subsection{Onsite Variation}

Varying the onsite strength shows that the effect persists over a range of parameters, with optimal agreement near:
\[
\text{onsite} \approx 1.0.
\]

\subsection{Two-Parameter Scan}

A scan over coupling $g$ and energy $E$ reveals a resonance line:
\[
E_\alpha(g),
\]
along which $\alpha^{-1}(T)$ remains stable.

\section{Interpretation}

Possible interpretations include:

\begin{itemize}
\item Fixed-point behavior in transport systems
\item Universality class of structured disorder
\item Hidden arithmetic constraints in lattice construction
\end{itemize}

At present, no theoretical explanation is known.

\section{Limitations}

\begin{itemize}
\item No derivation of $\alpha$ is provided
\item Mapping $T \mapsto \alpha^{-1}$ is model-dependent
\item Finite-size effects are not fully controlled
\item Dependence on prime set not fully explored
\end{itemize}

\section{Conclusion}

We report a robust numerical correlation between transport resonances and a dimensionless constant. While intriguing, the result remains empirical and requires further investigation.

\section*{Acknowledgments}

The author thanks exploratory computational experiments for motivating this work.

\end{document}